% Source: https://artofproblemsolving.com/community/c3h2460774_a_cool_geometry_problem
% Stable identity: aops-topic-2460774-q01
Let $P$, $Q$ be the centres of the circles of radius $5$ and $R$ the centre of the circle of radius $8$. Then $PQ=10$ and $PR=QR=13$. Place $P=(-5,0)$, $Q=(5,0)$; then $R$ is on the $y$-axis at height $\sqrt{13^2-5^2}=12$, so $R=(0,12)$.

Let the big circle have centre $Z$ and radius $\rho$. Internal tangency means
\[ZP=\rho-5,\qquad ZQ=\rho-5,\qquad ZR=\rho-8.\]
From $ZP=ZQ$, $Z$ lies on the perpendicular bisector of $PQ$, the $y$-axis: $Z=(0,k)$.
\begin{itemize}
\item If $k\le12$, then $ZR=12-k=\rho-8$, so $\rho=20-k$, and $ZP=\sqrt{25+k^2}=15-k$. Squaring, $25+k^2=225-30k+k^2$, so $k=\dfrac{20}{3}$ and $\rho=\dfrac{40}{3}$.
\item If $k>12$, then $k-12=\rho-8$, so $\rho=k-4$ and $\sqrt{25+k^2}=k-9$, which gives $25=-18k+81$, $k=\dfrac{28}{9}<12$: a contradiction.
\end{itemize}
Check: $ZP=\sqrt{25+\dfrac{400}{9}}=\dfrac{25}{3}=\rho-5$ and $ZR=12-\dfrac{20}{3}=\dfrac{16}{3}=\rho-8$. So the radius is $\boxed{\dfrac{40}{3}}$.

\begin{asy}
usepackage("amsmath");
size(5.5cm);
defaultpen(fontsize(10pt));
pair P=(-5,0), Q=(5,0), R=(0,12), Z=(0,20/3);
real rho=40/3;
draw(circle(Z,rho),red+linewidth(0.9));
draw(circle(P,5)); draw(circle(Q,5)); draw(circle(R,8));
dot(P); dot(Q); dot(R); dot(Z);
draw(P--Q,dashed); draw(P--R,dashed); draw(Q--R,dashed); draw(Z--P,blue);
draw(Z--(Z+rho*unit(Q-Z)),blue);
label("$P$",P,S); label("$Q$",Q,S); label("$R$",R,E); label("$Z$",Z,E);
label("$5$",(P+(-10,0))/2,N); label("$8$",(R+(0,20))/2,E);
draw(P--(-10,0)); draw(R--(0,20));
label("$\rho$",Z+0.75*rho*unit(Q-Z),NE,blue);
\end{asy}

(Descartes' circle theorem confirms this: with curvatures $\dfrac15,\dfrac15,\dfrac18$, the fourth curvature is $\dfrac{21}{40}\pm2\sqrt{\dfrac{1}{25}+\dfrac{1}{40}+\dfrac{1}{40}}=\dfrac{21}{40}\pm\dfrac35$; the value $-\dfrac{3}{40}$ belongs to the enclosing circle of radius $\dfrac{40}{3}$, and $\dfrac{9}{8}$ to the small circle of radius $\dfrac89$ in the gap.)
